%%%%%%%%%%%%%%%%%%%%%%%%%%%%%%%%%%%%%%%%%%%%%%%%%
\section{Evaluation Metrics for Fine-Grained Wildlife Detection}\label{sec:lit_evaluation_metrics}

Reporting a single mean-average-precision number is sufficient for a generic 80-class benchmark but is a poor fit for a 225-class, visually confusable, dual-domain wildlife taxonomy: a model that mistakes a plains zebra for a Grevy's zebra has made a fundamentally different error than one that misses the animal entirely, and a model evaluated on a synthetic image has been probed under a different data-generating process than one evaluated on a field photograph. This section surveys the literature underlying the thesis's evaluation design, moving from the general-purpose detection metrics that any object detector is judged against, through the diagnostic instruments developed for understanding \textit{why} a detector fails, to the question of how a test set that mixes real and synthetic imagery should be interpreted at all. The final subsection is the direct literature anchor for the evaluation methodology adopted in \Cref{sec:evaluation_framework} and summarized in \texttt{docs/plans/2026-06-10\_model-evaluation-strategy.md}.

\subsection{General Object Detection Metrics}\label{sec:lit_detection_metrics}

The dominant evaluation protocol for object detection was established by the \textit{PASCAL VOC} challenge \cite{everinghamPascalVisualObject2010}. A detection is judged correct if its bounding box overlaps a ground-truth box above a fixed \textit{Intersection-over-Union} (IoU) threshold -- the ratio of the intersection to the union of the two boxes -- with the threshold conventionally set to $0.5$, a deliberately low value chosen to absorb the inherent subjectivity of drawing a tight box around a non-convex object such as an animal with limbs extended \cite{everinghamPascalVisualObject2010}. Every prediction above the confidence threshold is then classified as a true or false positive, from which precision (the fraction of predictions that are correct) and recall (the fraction of ground-truth objects that are detected) are computed as functions of the confidence threshold; \textit{Average Precision} (AP) summarizes the resulting precision--recall curve as the mean of the interpolated precision at a fixed set of recall levels, and mean Average Precision (mAP) averages AP across classes \cite{everinghamPascalVisualObject2010}. This precision/recall/AP formulation was adopted specifically because a simple accuracy count is misleading under the heavy class imbalance between the small number of true objects and the very large number of candidate (non-object) regions a detector must reject \cite{everinghamPascalVisualObject2010}, and it has remained the basis of essentially every detection benchmark since \cite{zouObjectDetection202023}.

Two developments refined this baseline into the metric vector used throughout modern detection research. First, the fixed IoU$=0.5$ threshold rewards coarse localization and does not distinguish a tightly fitted box from a loosely fitted one, so long as both clear the 50\% bar; the \textit{COCO} benchmark \cite{linMicrosoftCOCOCommon2015} addressed this by averaging AP over ten IoU thresholds spanning $0.5$ to $0.95$ in steps of $0.05$ -- denoted AP@[.5:.95] -- so that a detector is rewarded for progressively tighter localization rather than merely clearing a single bar \cite{zouObjectDetection202023}. The single-threshold AP50 and the stricter AP75 are retained alongside the averaged AP@[.5:.95] as complementary operating points, the former closer to the original PASCAL protocol and the latter probing high-precision localization specifically \cite{zouObjectDetection202023}. Second, because AP only credits a detector for the boxes it does emit, it says nothing about how many ground-truth objects a detector never proposes at all irrespective of confidence; \textit{Average Recall} (AR), typically reported at a fixed budget of the top $1$, $10$, or $100$ detections per image (AR@1/10/100), complements AP by measuring this upper bound on achievable recall independent of the confidence ranking \cite{zouObjectDetection202023}. Together, the COCO-style twelve-statistic vector -- AP@[.5:.95], AP50, AP75, AP split by object size (small/medium/large), and AR@1/10/100 split the same way -- is the de facto standard instrument against which any new detector, including the \textit{YOLO}-family models surveyed in \Cref{sec:lit_yolo_family}, is reported \cite{zouObjectDetection202023}. This full vector is exactly the instrument this thesis retains as its own baseline reporting format (\Cref{sec:evaluation_framework}), precisely because it is the bridge that keeps the thesis's numbers legible against the published literature -- while, as the remaining subsections argue, being too coarse on its own to answer the questions a 225-class look-alike-heavy wildlife taxonomy actually raises.

\subsection{Hierarchical and Coarse-to-Fine Evaluation for Confusable Classes}\label{sec:lit_hierarchical_evaluation}

A single flat AP number computed over all 225 classes cannot distinguish a model that finds an animal and names the wrong species within a visually confusable group (e.g.\ a plains zebra scored as a Grevy's zebra) from one that names a completely unrelated class (a wolf scored as a squirrel) -- yet the two errors are, from a practical-deployment standpoint, not remotely equivalent. The literature on diagnosing detector error provides an early and directly applicable precedent for this distinction, even though it predates today's large-scale fine-grained taxonomies. \textcite{hoiemDiagnosingErrorObject2012}, analyzing false positives on \textit{PASCAL VOC}, explicitly separate ``confusion with similar objects'' from other false-positive types by defining, \textit{a priori} and independently of any single detector's actual mistakes, small sets of semantically related categories (their example groupings include \{all vehicles\}, \{all animals including person\}, and \{chair, diningtable, sofa\}); a false positive falling within the same set as its ground-truth neighbor is counted separately from one that is not. Their results show that this within-group confusion, together with pure localization error, accounts for the large majority of high-confidence false positives -- markedly more than confusion with dissimilar, unrelated background clutter -- which is precisely the failure mode this thesis expects to dominate among the zebra, elephant, lynx-caracal, and gazelle/antelope look-alike clusters in the target taxonomy.

The critical methodological feature of Hoiem et al.'s grouping is that it is defined \textit{before} any detector is run, from domain knowledge about which categories are visually or semantically related, rather than being read off a trained model's own confusion matrix after the fact. This is precisely the distinction this thesis's own grouping strategy must make explicit when defining its look-alike groups: a taxonomic or curated-visual-similarity rollup, decided independently of any candidate model's behavior, versus an empirically-derived grouping mined from a model's confusion matrix, which would be circular if used to define the very categories that model is then evaluated against. The literature surveyed for this thesis does not contain a paper that performs this rollup at species/genus/family taxonomic resolution specifically, or that formally compares a taxonomy-driven grouping against an empirical-confusion-driven one as competing evaluation designs -- the closest available anchor is the general a-priori-grouping principle established by \textcite{hoiemDiagnosingErrorObject2012}, which this thesis extends from a handful of hand-picked semantic sets to a full taxonomic hierarchy (\Cref{sec:evaluation_framework}). This is stated plainly here rather than papered over with a tangentially related citation.

\subsection{Error-Type Decomposition}\label{sec:lit_error_decomposition}

A single AP figure is a compressed summary that discards the information needed to say \textit{why} a detector under-performs -- whether it fails to localize objects at all, localizes them but assigns the wrong label, or duplicates detections on the same object. \textcite{hoiemDiagnosingErrorObject2012} were among the first to formalize this decomposition for object detection, partitioning every false positive into one of four mutually exclusive categories: localization error (a detection that overlaps a true object of the correct class but with insufficient IoU, which they treat as including duplicate detections on an already-detected object), confusion with a semantically similar object, confusion with a dissimilar labeled object, or confusion with unlabeled background. Reporting the frequency of each category, and separately the AP improvement that would follow from eliminating each one in isolation, lets a researcher identify which failure mode is most worth addressing rather than relying on a single aggregate number -- their own analysis, for instance, shows that correcting localization error and similar-object confusion would yield far larger AP gains than eliminating background false positives entirely, a finding that would be invisible in mAP alone.

This four-way error taxonomy is the direct conceptual precedent for the automated, off-the-shelf error-decomposition toolboxes now used routinely in detection research, which extend the same classification/localization/duplicate/background partition (with an added ``missed ground truth'' category) into a single reusable software pipeline applicable to any detector's raw predictions without the manual re-annotation Hoiem et al.\ required. This thesis's own evaluation pipeline (\Cref{sec:evaluation_framework}) adopts exactly this kind of decomposition as a complement to the granularity-gap analysis of \Cref{sec:lit_hierarchical_evaluation}, using it to corroborate whether an accuracy loss between the coarse and fine granularity levels is actually attributable to classification confusion among look-alike species, as expected, rather than to a hidden localization or duplicate-detection problem. The literature survey conducted for this thesis was not able to establish a verifiable, corpus-confirmed citation for the specific modern toolbox implementation used operationally in \Cref{chapter:results}; this subsection therefore anchors the underlying method in \textcite{hoiemDiagnosingErrorObject2012} rather than name a downstream implementation without a matching source.

\subsection{Evaluating Across a Real/Synthetic Test-Domain Split}\label{sec:lit_real_synthetic_evaluation}

Every evaluation protocol discussed so far, from \textit{PASCAL VOC} to \textit{COCO} to Hoiem et al.'s error diagnostics, is computed over a single, homogeneous test set of real photographs. This thesis's dataset instead ships two test domains -- an uneven, per-class real test set and a class-balanced $225\times50$ synthetic test set (\Cref{sec:synthetic_data_supplementation}) -- and \Cref{sec:evaluation_framework} adopts the union of the two, scored together as one set, as the thesis's default headline evaluation condition, with the real-only subset always reported alongside it as the primary-evaluation figure and the anchor for any comparison to public, real-image benchmarks. Before that design is used, it is worth being explicit about what precedent for it does, and does not, exist in the surveyed literature.

Standard practice when synthetic imagery is involved in training is to keep evaluation strictly on real, established test sets rather than to blend synthetic images into the test condition itself. Domain-randomization work aimed at exactly the same sim-to-real problem this thesis faces -- training a detector on synthetic imagery and hoping it generalizes to the real deployment domain -- evaluates its domain-randomized detector exclusively on the real-world \textit{KITTI} test set, never mixing synthetic renders into the test condition, even though a synthetic dataset rendered to match the test domain (\textit{Virtual KITTI}) was available as a comparison point \cite{tremblayTrainingDeepNetworks2018}. Similarly, work evaluating whether modern generative-model output is usable as training data for image recognition benchmarks its synthetically pretrained or synthetically trained models against standard, unmodified real test sets (\textit{CIFAR-100}, \textit{PASCAL VOC}) rather than any blended variant, and explicitly attributes synthetic-to-real performance gaps to a domain gap between the generated and real image distributions \cite{heSyntheticDataGenerative2023}. The general problem of a mismatch between the distribution a model is evaluated on and the distribution it will actually encounter -- of which a train/test domain shift between synthetic and real imagery is one instance -- is itself the subject of a substantial domain-generalization literature \cite{zhouDomainGeneralizationSurvey2023} and traces back to early work establishing that benchmark datasets themselves carry systematic biases that limit how their measured performance transfers to other distributions \cite{UnbiasedLookDataset}. None of this literature, however, treats a synthetic partition as a component to be folded into a single headline test metric alongside real data; where synthetic data appears in a test protocol at all, it is either the sole test domain (to isolate a specific capability) or is kept rigidly separate from the real test condition precisely so that a domain-shift gap remains measurable.

Taken together, the surveyed literature offers \textbf{no established precedent} for two specific choices this thesis makes: reporting a headline detection metric computed over the union of a real and a synthetic test set, and using a class-balanced synthetic test set as a diagnostic instrument for classes where the real test data is thin or noisy. Both are presented here, deliberately, as explicitly justified methodological choices of this thesis rather than as applications of an established practice. The justification itself -- that a fixed, class-balanced synthetic probe stabilizes an otherwise high-variance per-class metric for the data-poor Band~A/B species while remaining a negligible addition for the well-resourced Band~D species, and that the real-only breakout is always reported alongside the mixed figure specifically so this choice cannot silently inflate the headline number -- is developed in full in \Cref{sec:evaluation_framework} and \texttt{docs/plans/2026-06-10\_model-evaluation-strategy.md}. The adjacent literature surveyed above is used here only as context for that design decision -- to show what the established alternative (real-only evaluation, kept strictly separate from any synthetic probe) looks like, and why departing from it needs its own argument -- and not as literature that itself endorses a mixed real+synthetic headline metric. Consistent with this, \Cref{sec:evaluation_framework} treats the real-vs-synthetic performance delta as a standing watchdog: should it reveal the synthetic component to be systematically inflating or deflating the mixed headline relative to real-world performance, the default evaluation axes are to be revised.
